\section{Encaminamiento multiclase: las coaliciones son cuerpos, no puntos}
\label{sec:multiclase}

El tratamiento del tráfico en \cref{sec:megajuego} penaliza la congestión de
forma agregada, lo que basta cuando todos los agentes ocupan lo mismo. No es el
caso aquí: una coalición en modo \textsf{cargo} alrededor de una carga ancha
ocupa un corredor que un AMR libre apenas roza. Ignorar esa diferencia hace que
el mecanismo cobre lo mismo por consumos de capacidad muy distintos.

\subsection{Clases de envolvente}

\begin{definicion}[Clase y equivalente de capacidad]
\label{def:clase}
Una coalición que transporta la carga $\ell$ con su formación define una
envolvente de ancho $w^{\mathrm{env}}_\ell$. Con $w_0$ el ancho de un AMR libre,
su \emph{clase} es $c=\lceil w^{\mathrm{env}}_\ell/w_0\rceil$ y su equivalente de
capacidad
\[
  p_c=\frac{w^{\mathrm{env}}_c}{w_0}\ \ge\ 1 .
\]
El flujo de capacidad del corredor $e$ es $y_e=\sum_c p_c\,x^c_e$, donde $x^c_e$
es el flujo de clase $c$. El corredor admite la clase $c$ solo si
$W_e\ge w^{\mathrm{env}}_c$, de modo que cada clase circula sobre su propio
subgrafo de transitabilidad $\Gcal^c$.
\end{definicion}

El tiempo de traversa depende del flujo de capacidad, no del número de cuerpos:
$t_e(y_e)=t^0_e\bigl(1+\alpha(y_e/C_e)^{\beta}\bigr)$.

\subsection{Peaje marginal por clase}

\begin{supuesto}[Coste generalizado]
\label{s:coste-clase}
El coste que percibe una unidad de flujo de clase $c$ al recorrer $e$ es
$p_c\,t_e(y_e)$: proporcional al tiempo y a la capacidad que consume.
\end{supuesto}

\begin{teorema}[El peaje marginal por clase alinea el equilibrio con el óptimo]
\label{th:multiclase}
Bajo \cref{s:coste-clase}, con peajes
\[
  \tau^c_e\ =\ p_c\,y_e\,t'_e(y_e),
\]
todo óptimo social interior del problema
$\min\sum_e y_e\,t_e(y_e)$ satisface la condición de Wardrop para cada clase:
el coste con peaje $p_c t_e+\tau^c_e$ está igualado sobre los corredores que la
clase $c$ utiliza. El encaminamiento egoísta de cada coalición sobre $\Gcal^c$
es, por tanto, socialmente óptimo.
\end{teorema}

\begin{proof}
La derivada del coste social respecto a $x^c_e$ es
$\partial_{x^c_e}\sum_{e'}y_{e'}t_{e'}(y_{e'})
=p_c\bigl[t_e(y_e)+y_e t'_e(y_e)\bigr]$.
En un óptimo interior esa derivada está igualada sobre los corredores usados por
la clase $c$, luego el corchete $\bigl[t_e+y_et'_e\bigr]$ lo está. El coste con
peaje es $p_c t_e+p_c y_e t'_e=p_c\bigl[t_e+y_e t'_e\bigr]$, es decir $p_c$ por
ese mismo corchete, y queda igualado para toda clase simultáneamente.
\end{proof}

\begin{observacion}[Comprobación numérica]
\label{obs:multiclase-num}
Con dos corredores en paralelo, función BPR
($\alpha=0{,}15$, $\beta=4$, $t^0=(1{,}0,\,1{,}4)$, $C=(3{,}0,\,4{,}0)$) y dos
clases de equivalentes $p=(1{,}0,\,2{,}5)$ con demandas $(2{,}0,\,1{,}0)$, el
óptimo social reparte $y=(2{,}639,\,1{,}861)$ y el corchete
$t_e+y_et'_e$ vale $1{,}449$ en ambos corredores. El coste con peaje queda
igualado en $1{,}449$ para la clase libre y en $3{,}623$ para la clase ancha.
\end{observacion}

\begin{observacion}[Qué es clásico y qué no]
\label{obs:multiclase-atrib}
La asignación de tráfico multiclase con equivalentes de capacidad y peaje
marginal es teoría establecida desde \cite{dafermos1972multiclass}; el
\cref{th:multiclase} es su transcripción. Lo que corresponde a este trabajo es
la interpretación de la clase: no un tipo de vehículo, sino la
\emph{envolvente de una coalición cooperativa}, que además es endógena —cambia
con el modo, con la geometría de contactos y con la reconfiguración de puestos—
y observable localmente, de modo que $x^c_e$ se estima con el mismo consenso de
\cref{th:mapa}. Y el acoplamiento con el reclutamiento: $t_e$ depende de la
velocidad de la coalición, que es la de su miembro más lento por
\cref{le:curvatura}, con lo que la elección de ruta y la de miembros dejan de
ser separables.
\end{observacion}

\subsection{Admisión de tareas como estabilizador de precios}

\Cref{pr:dominancia} da la condición bajo la cual la red de precios es estable,
pero no dice cómo mantenerla. Hay una regla operativa directa.

\begin{corolario}[Control de admisión]
\label{co:admision}
Si el lanzamiento del reclutamiento de una carga $\ell$ se autoriza solo cuando
\[
  E_{\ell\ell}\ \ge\ \sum_{\ell'\ \mathrm{activas}}\bigl|E_{\ell\ell'}\bigr|+m,
  \qquad m>0,
\]
y las cargas que no lo cumplen esperan en una cola de admisión, entonces la
matriz de elasticidades permanece diagonalmente dominante en todo instante y,
por \cref{pr:dominancia}, el ajuste de precios es estable.
\end{corolario}

El precio de esta regla es explícito y hay que declararlo: la espera en cola se
suma al plazo de la carga y compite con su urgencia. La estabilidad de precios
se compra con latencia de lanzamiento.

\begin{observacion}[Corrección sobre el certificado de la dinámica de revisión]
\label{obs:smith-acoplado}
\Cref{obs:smith} justifica el uso de Smith por estacionariedad de Nash y por
satisfacer la hipótesis de los teoremas de convergencia al GNE. Conviene
precisar el alcance: esas condiciones se establecen para juegos sin
restricciones acopladas. Con el conjunto admisible $\Fcal(t)$ acoplado —que es
el caso aquí— el certificado aplicable no es el de la dinámica de revisión
aislada, sino el de su composición con la dinámica de multiplicadores
distribuidos que gestiona las restricciones compartidas
\cite{martinezPiazuelo2025gne}. El argumento de \cref{th:precio-seguridad}
—multiplicador común por restricción compartida— es precisamente lo que hace
lícita esa composición.
\end{observacion}

\subsection{El conflicto de corredor como negociación}

La terna lexicográfica que resuelve un cruce es una regla de prioridad, y toda
regla de prioridad es un punto extremo del conjunto de acuerdos posibles.
Conviene ver dónde queda respecto de los criterios de negociación, porque el
conflicto de corredor tiene exactamente la estructura de un problema de regateo:
dos partes, ganancias distintas según quién pase primero, y un \emph{punto de
desacuerdo} que es el bloqueo mutuo \cite{clempner2024optimization}.

\begin{observacion}[Tres formas de resolver el mismo cruce]
\label{obs:corredor-negociacion}
Con dos coaliciones que obtienen $10$ y $6$ al pasar primero, $4$ y $3$ al
esperar, y $(1;\ 0{,}5)$ en caso de bloqueo:
\[
  \begin{array}{lccc}
    \text{criterio} & \text{pagos} & \text{suma} & \text{producto de excedentes}\\\hline
    \text{lexicográfico (mecanismo actual)} & (10;\ 3) & 13{,}00 & 22{,}50\\
    \text{Nash} & (8;\ 4) & 12{,}00 & 24{,}50\\
    \text{Kalai--Smorodinsky} & (7{,}30;\ 4{,}35) & 11{,}65 & 24{,}26
  \end{array}
\]
La regla de prioridad es la que \emph{más caudal total produce} —$13$ frente a
$12$— y a la vez la más desigual. Los criterios de negociación compran equilibrio
al precio de un $8\%$ de caudal.
\end{observacion}

\begin{observacion}[Cómo se implementa un punto interior]
\label{obs:alternancia}
El punto de Nash no corresponde a ninguna decisión pura: está en el interior del
segmento entre «pasa la primera» y «pasa la segunda», y se realiza
\emph{aleatorizando} quién cede. En un almacén, donde los mismos pares de rutas
se cruzan muchas veces al día, eso no exige moneda: se implementa
\emph{alternando} con la frecuencia adecuada sobre encuentros repetidos.

Esto da una lectura útil de la terna lexicográfica. Su desempate por
identidad —el último criterio, $-i$— es determinista y por tanto favorece
sistemáticamente al mismo AMR en empates. Sustituirlo por una alternancia con la
proporción de \cref{obs:corredor-negociacion} convertiría el mecanismo en una
implementación repetida de la solución de negociación, a coste nulo en
comunicación. Es una modificación pequeña con consecuencia distributiva, y se
declara como opción de diseño, no como mejora demostrada de caudal: el caudal,
de hecho, baja.
\end{observacion}

\subsection{De la geometría al número de rondas}

La reserva de corredores es, vista como problema combinatorio, una coloración: dos
coaliciones cuyos trayectos se solapan no pueden tener el mismo turno, y el número
mínimo de turnos libres de conflicto es el número cromático del grafo de conflicto
\cite{yadav2023advanced}. Calcularlo exactamente es NP-completo, pero no hace
falta.

\begin{proposicion}[Cota geométrica del número de turnos]
\label{pr:rondas-color}
Sea $\Delta$ el grado máximo del grafo de conflicto. La coloración voraz por
grados decrecientes usa a lo sumo $\Delta+1$ turnos, y por tanto
$\chi\le\Delta+1$. Componiendo con la cota de grado de
\cref{pr:compatibilidad}, el número de turnos necesarios está acotado por una
cantidad puramente geométrica:
\[
  \#\text{turnos}\ \le\ 1+\Bigl\lfloor \pi\big/\arcsin\bigl(d_{\min}/(2R^{\mathrm{com}})\bigr)\Bigr\rfloor .
\]
\end{proposicion}

\begin{observacion}[Los números, y la comprobación]
\label{obs:rondas-num}
Para razones $d_{\min}/R^{\mathrm{com}}$ de $0{,}05$ a $0{,}5$ la cota da $126$,
$63$, $32$, $21$ y $13$ turnos. Sobre trescientos grafos de conflicto generados
aleatoriamente, la coloración voraz nunca excedió $\Delta+1$.

Lo relevante no es la cota en sí —es holgada— sino que sea \emph{geométrica}: el
número de turnos que un corredor puede necesitar queda determinado por la
separación mínima y el alcance de comunicación, dos parámetros de diseño de la
instalación. Junto con \cref{obs:ciclo-despeje}, que fija la duración de cada
turno, se obtiene una cota del tiempo de espera en un cruce a partir de
magnitudes que se conocen antes de desplegar la flota.
\end{observacion}

\subsection{Por qué las revisiones individuales no bastan}

Hay una segunda lectura combinatoria, esta sobre la asignación misma. Emparejar
AMR con cargas es un problema de emparejamiento, y el teorema de Berge lo
caracteriza: un emparejamiento es máximo si y solo si no admite \emph{camino
aumentante} \cite{yadav2023advanced}. Un camino aumentante es una cadena de
reasignaciones, no un cambio individual.

\begin{observacion}[Estable ante desviaciones individuales y sin embargo no máximo]
\label{obs:berge}
El caso mínimo lo muestra sin margen de duda. Sean dos AMR y dos cargas, donde
el primero puede servir ambas y el segundo solo la primera. Con el primero
asignado a la carga que ambos pueden servir, ningún AMR libre puede tomar una
carga libre: la asignación es estable ante revisiones individuales. Su tamaño es
$1$, y el máximo es $2$ —basta que el primero ceda esa carga y tome la otra—.
Una carga queda sin servir pudiendo servirse.

Esto es la versión combinatoria de \cref{obs:fuerte-ciervo}: allí el mal
equilibrio resistía desviaciones individuales y caía ante una coalición; aquí la
asignación subóptima resiste revisiones individuales y cae ante una cadena. Son
el mismo hecho en dos lenguajes, y ambos apuntan al mismo remedio —permitir
desviaciones coordinadas— que en esta arquitectura proporciona la guarda de
\cref{obs:guarda-fuerte}.

De paso acota lo que \cref{co:terminacion} garantiza: termina, y termina en un
punto estable ante desviaciones de la clase que el protocolo admite. Ampliar esa
clase a cadenas de longitud acotada es una mejora identificada y no implementada.
\end{observacion}

\subsection{Cuánto debe durar la reserva}

La asignación de derecho de paso entre flujos que compiten por un mismo espacio
es el problema de sincronización de semáforos, cuyo objetivo es minimizar la cola
total \cite{clempner2024optimization}. Trasladarlo aquí obliga a incorporar algo
que el caso vehicular puede despreciar y este no.

\begin{observacion}[El despeje fija la duración del turno]
\label{obs:ciclo-despeje}
Un vehículo abandona el cruce en un instante; una coalición ancha y lenta
necesita un tiempo de \emph{despeje} $L$ para salir del corredor, y durante ese
tiempo nadie más puede entrar. Ese coste se paga en cada conmutación. Minimizando
el retardo total sobre la duración del ciclo $C$ con reparto de verde
proporcional a la demanda, el óptimo resulta
\[
  C^\star\ \approx\ 5{,}9\,L,
\]
con una constancia notable en $L$: la razón se mantiene entre $5{,}9$ y $6{,}0$ al
variar $L$ desde $0{,}5$ hasta $30$~s, sobre un rango de sesenta veces. No es
constante en la \emph{saturación} $Y=\sum_i y_i$: un barrido posterior da
$C^\star/L\approx4{,}33/(1-Y)-0{,}57$ ($R^2=0{,}998$), es decir $4{,}7$, $5{,}9$,
$8{,}3$ y $20{,}7$ para $Y=0{,}2$, $0{,}32$, $0{,}5$ y $0{,}8$; el valor $5{,}9$
corresponde a $Y=0{,}32$. El
retardo mínimo alcanzable crece linealmente con $L$: de $0{,}391$ a $23{,}46$ en
ese mismo barrido.

De aquí sale una regla de dimensionamiento que el mecanismo no tenía: \emph{el
testigo de reserva debe retenerse del orden de seis veces el tiempo de despeje de
la coalición más ancha que pueda usar el corredor}. Y justifica la elección de
un testigo exclusivo frente a una alternancia rápida: con despeje grande el ciclo
óptimo es largo, y un ciclo largo con reparto proporcional \emph{es} una
asignación casi exclusiva por turnos. La alternancia rápida solo conviene cuando
el despeje es pequeño, que es el caso vehicular y no el de una coalición
acoplada.
\end{observacion}

\subsection{Exclusión, vivacidad e interbloqueo}

El testigo de corredor garantiza exclusión. Conviene decir con precisión qué
más garantiza y qué no \cite{mayorga2026compendio}.

\begin{proposicion}[Invariante de un arrendamiento versionado]
\label{pr:arrendamiento}
Sea un recurso $z$ con estado $(o_z,\nu_z,t_z^{\exp})$: propietario o vacío,
versión y expiración. Si (i) una concesión solo se confirma cuando el recurso
está libre o su arrendamiento anterior expiró, (ii) cada \emph{commit}
incrementa estrictamente $\nu_z$, y (iii) los mensajes con versión antigua se
descartan, entonces en todo estado alcanzable hay a lo sumo un propietario
vigente.
\end{proposicion}

\begin{proof}
Inducción sobre \emph{commits}. Una liberación deja cero propietarios; una
renovación conserva el mismo; una nueva concesión solo se confirma tras
observar libre o expirado y genera una versión nueva, de modo que cualquier
concesión concurrente basada en la versión anterior queda obsoleta y se
descarta.
\end{proof}

\begin{proposicion}[Exclusión no implica vivacidad]
\label{pr:exclusion-no-vivacidad}
Existe una ejecución que satisface \cref{pr:arrendamiento} para todos los
recursos y permanece bloqueada indefinidamente: dos recursos unitarios, la
coalición $1$ posee $z_1$ y espera $z_2$, la coalición $2$ posee $z_2$ y
espera $z_1$. Ningún recurso tiene dos propietarios y ninguna coalición
progresa.
\end{proposicion}

\begin{teorema}[Ciclo de espera e interbloqueo]
\label{th:ciclo-espera}
Con recursos de instancia única, retención mientras se espera, ausencia de
expropiación y liberación solo al completar la operación esperada, un
conjunto cerrado y persistentemente bloqueado de coaliciones contiene un ciclo
en el grafo de espera. Recíprocamente, un ciclo dirigido persistente basta
para el interbloqueo de las coaliciones del ciclo.
\end{teorema}

\begin{proof}
En un conjunto cerrado y bloqueado cada coalición espera un recurso poseído
por otra del mismo conjunto, luego cada nodo tiene una arista saliente en un
grafo finito y seguir aristas repite un nodo. En un ciclo
$c_1\to\cdots\to c_m\to c_1$, cada una espera un recurso retenido por la
siguiente; sin expropiación ni liberación en espera, ninguna completa la
operación que liberaría, y el ciclo persiste.
\end{proof}

\begin{observacion}[Equidad y ausencia de interbloqueo son contratos distintos]
\label{obs:equidad-interbloqueo}
Una regla de envejecimiento elimina la inanición en la cola de un recurso y no
rompe un ciclo de \cref{th:ciclo-espera} que atraviese dos recursos. Para el
corredor, la consecuencia de diseño es concreta: el testigo debe adquirirse
en un orden total de recursos, o liberarse mientras se espera —lo que es la
alternancia de \cref{obs:alternancia}—, porque la exclusión por sí sola no
garantiza que la carga llegue.
\end{observacion}

\subsection{Potencial exacto con pesos y latencias afines}

\begin{proposicion}[Juego atómico ponderado afín]
\label{pr:potencial-afin-ponderado}
Sean coaliciones atómicas $i=1,\dots,N$ con peso físico $w_i>0$ y estrategia
$s_i\subseteq E$ (los recursos que ocupa), carga $y_e(s)=\sum_{i:e\in s_i}w_i$
y latencia afín $\ell_e(y)=a_ey+b_e$ con $a_e,b_e\ge0$. Con coste
$C_i(s)=w_i\sum_{e\in s_i}\ell_e(y_e(s))$, el juego tiene potencial exacto
\begin{equation}
  \label{eq:potencial-afin}
  \Phi(s)=\sum_{e\in E}\Bigl[\frac{a_e}{2}\Bigl(y_e(s)^2+\sum_{i:e\in s_i}w_i^2\Bigr)+b_ey_e(s)\Bigr],
\end{equation}
y posee al menos un equilibrio puro al que toda sucesión de mejoras
unilaterales estrictas llega en tiempo finito.
\end{proposicion}

\begin{proof}
Basta un recurso. Si $i$ entra a $e$ con carga previa $y$, el potencial
crece $\tfrac{a_e}{2}[(y+w_i)^2-y^2+w_i^2]+b_ew_i=a_eyw_i+a_ew_i^2+b_ew_i
=w_i\ell_e(y+w_i)$, que es su coste en $e$; si sale, la variación es el
negativo de su coste actual. En mil desviaciones aleatorias el error de la
identidad es $1{,}8\times10^{-14}$. Existencia y mejora finita siguen de la
finitud de $E$.
\end{proof}

\begin{observacion}[Dónde acaba esta clase]
\label{obs:afin-limite}
Con pesos heterogéneos el potencial exacto existe para latencias afines y
no en general: es la frontera clásica de los juegos de congestión ponderados
\cite{fotakis2005selfish}. Para latencias no lineales, la vía exacta es la
externalidad atómica de \cref{th:externalidad-atomica}, que cobra la
contribución marginal en lugar de la latencia; el precio es que el coste ya
no es «lo que tardo» sino «lo que causo».
\end{observacion}

\subsection{Huellas extensas: de los centros a los conjuntos}

Una coalición no es un punto. Lo que sigue fija cuándo la abstracción puntual
es válida y cuánto cuesta cuando no lo es \cite{mayorga2026compendio}.

\begin{lema}[Obstáculo de configuración y obstáculo relativo]
\label{le:minkowski}
Fijada la orientación, un centro $p$ es libre para la huella $\bar F$ si y
solo si $p\notin O\oplus(-\bar F)$; dos huellas con centros $p_c,p_d$
colisionan si y solo si $p_c-p_d\in\bar F_d\oplus(-\bar F_c)$, que para
discos de radios $R_c,R_d$ es $\lVert p_c-p_d\rVert\le R_c+R_d$. En
consecuencia, si $\bar F_1\subseteq\bar F_2$ entonces
$Q_{\mathrm{free}}(\bar F_2)\subseteq Q_{\mathrm{free}}(\bar F_1)$: agrandar
una huella nunca crea caminos.
\end{lema}

\begin{proof}
Hay colisión si y solo si existen $o\in O$ y $f\in\bar F$ con $o=p+f$, es
decir $p\in O\oplus(-\bar F)$; la versión relativa es $p_c+f_c=p_d+f_d$; la
suma de Minkowski es monótona por inclusión y tomar complementos la invierte.
\end{proof}

\begin{proposicion}[Certificado de despeje y su contraejemplo mínimo]
\label{pr:certificado-despeje}
Si $\bar F_c\subseteq B_{R_c}(0)$ y la curva de centros cumple
$\mathrm{dist}(p(t),O)>R_c+d_{\mathrm{safe}}$ para todo $t$, la huella
completa conserva margen $d_{\mathrm{safe}}$. La condición no es vacía en el
otro sentido: para todo $R>0$ existe un corredor recto de ancho $0<w<2R$ en
el que el eje central es un camino puntual y ningún disco de radio $R$ lo
recorre.
\end{proposicion}

\begin{proof}
Para $z$ en la huella, $\mathrm{dist}(z,O)\ge\mathrm{dist}(p,O)-\lVert z-p\rVert>d_{\mathrm{safe}}$.
En el corredor, todo disco de radio $R$ toca una pared porque su diámetro
excede $w$, y por \cref{le:minkowski} el espacio de centros libres es vacío.
\end{proof}

\begin{observacion}[El despeje normalizado separa exactamente los falsos positivos]
\label{obs:despeje-normalizado}
Con $\zeta=d_{\min}(\pi,O)/(R_c+d_{\mathrm{safe}})$, en el banco de cinco
geometrías, cuarenta semillas y siete escalas de huella, todos los casos con
$\zeta\le1$ fueron falsos positivos de la representación puntual y no hubo
ninguno con $\zeta>1$; agregando radios, la tasa de falso positivo fue del
$13{,}2\%$, y del $39{,}5\%$ para la huella grande, mientras la búsqueda con
huella encontró camino válido en el $100\%$ y alargó el $9{,}5\%$ de los
pares en a lo sumo cuatro celdas \cite{mayorga2026compendio}. La búsqueda A*
sigue siendo exacta para el grafo que recibe; lo que falla es el grafo cuando
el despeje deja de dominar la huella de la coalición.
\end{observacion}

\begin{lema}[Cota de caminos independientes y su límite]
\label{le:caminos-independientes}
Si $d_c$ es la longitud del camino individual más corto del agente $c$
ignorando a los demás, toda solución conjunta factible $\pi$ cumple
$\mathrm{SOC}(\pi)\ge\sum_cd_c$; la igualdad no implica factibilidad
conjunta, y un plan sin coincidencia de centros sigue siendo inviable si en
algún instante $\lVert p_i(t)-p_j(t)\rVert<R_i+R_j$.
\end{lema}

\begin{proposicion}[La planificación por prioridades no es completa]
\label{pr:prioridad-incompleta}
Existe una instancia factible que un orden fijo de prioridades no resuelve:
un corredor en T con celdas $\{\ell,c,r,u\}$ y aristas $\ell-c-r$, $u-c$; el
agente $1$ va de $u$ a $c$ y el $2$ de $\ell$ a $r$. Con prioridad para $1$,
su camino $u\to c$ se reserva y al quedarse en su objetivo bloquea para
siempre el único paso $\ell\to c\to r$; la solución conjunta existe: $1$
espera en $u$, $2$ ejecuta $\ell\to c\to r$ y después $1$ entra.
\end{proposicion}

\begin{observacion}[Tres campañas, tres modos de fallo]
\label{obs:campanas-n123}
En ciento veinte mundos pareados con dos a cuatro coaliciones, la búsqueda
conjunta con detector de agentes extensos completó el $100\%$; la misma
búsqueda sobre centros con comprobación posterior de huella, el $66{,}7\%$;
la independiente, el $8{,}3\%$; y la corrección de agentes extensos produjo
cuarenta discordancias favorables y ninguna adversa, con McNemar exacto
$p=2^{-39}=1{,}8\times10^{-12}$. Con una implementación pública de
prioridad con herencia, el $100\%$ de éxito puntual cayó al $50\%$ al
auditar huellas heterogéneas. Y con información vecinal en sesenta mundos,
la evitación reactiva completó el $50\%$ con $35\%$ de colisiones y $15\%$ de
interbloqueos, mientras el arrendamiento local completó el $16{,}7\%$ con cero
colisiones y $83{,}3\%$ de interbloqueos \cite{mayorga2026compendio}. La
última fila es \cref{pr:exclusion-no-vivacidad} medida: la exclusión compra
seguridad y paga vivacidad, y la elección entre las dos filas no es de
algoritmo sino de contrato, que es lo que el testigo de corredor con orden
total de \cref{obs:equidad-interbloqueo} resuelve.
\end{observacion}
